% Original research note for the present investigation.  This fragment is
% self-contained modulo amsmath, amssymb, and amsthm.  It does not claim
% proof-assistant verification or literature-wide priority.

\section{A finite-model obstruction to eager initialization}
\label{sec:eager-obstruction}

There is a tempting way to shorten the computation clause of an iteration
checker.  Instead of preserving an independent initializer at counter zero,
apply the update immediately:
\begin{align}
 F(0,a;v.h)&=h[a/v],\label{eq:eager-zero}\\
 F(Sn,a;v.h)&=h[F(n,a;v.h)/v].\label{eq:eager-step}
\end{align}
The same syntactic opening operation then appears in both computation
clauses.  This subsection proves that, with the stated remaining
primitives, this change cannot present iPRA.  The obstruction applies to
relational interpretations and to finite-dimensional interpretations with
quotients, as well as to translations by terms.

\paragraph{The precise proposed theory.}
Let the numerical terms be the finite expressions
\[
 t ::= x\mid 0\mid S(t)\mid F(t,a;v.h).
\]
In the last constructor, $v$ is bound in the update $h$ only.  Every other
free variable of $h$ is a passive parameter.  Nested occurrences of $F$
are allowed, but this is a finite structural grammar: the update does not
refer circularly to the operation currently being defined.  Equivalently,
one may use a staged first-order signature containing a separate operation
symbol for each already constructed update description.  Substitution is
capture avoiding.

Let $T_{\mathrm e}$ use ordinary intuitionistic first-order logic with
actual equality, the equations \eqref{eq:eager-zero}--\eqref{eq:eager-step},
the two boundary assertions
\[
 \forall x\,(x=0\vee\exists y\,x=Sy),\qquad 0\ne S0,
\]
and induction on literal equations, with arbitrary retained assumptions
and passive parameters.  The countermodel below also satisfies the
stronger assertion $\forall x\,(0\ne Sx)$ and induction on \emph{every}
formula.  It therefore covers any version obtained by strengthening the
induction schema in that way.  Successor injectivity is \emph{not} an
axiom of this proposed simplification.

\paragraph{Two-element interpretation.}
Take the ordered set $M=\{0,1\}$ with $0<1$, interpret zero by $0$, and put
\[
                         S^M(b)=1\quad(b\in M).
\]
For a variable assignment $\rho$, interpret the eager constructor by
\[
 \llbracket F(n,a;v.h)\rrbracket_\rho
   :=\llbracket h\rrbracket_{\rho[v\mapsto\llbracket a\rrbracket_\rho]}.
 \tag{Eager interpretation}
\]
The counter is deliberately ignored.  This equation is a definition of
the mathematical interpretation, not a new evaluator trusted by the
proposed proof checker.

\paragraph{Monotonicity lemma.}
Every numerical term is monotone in all its free variables under this
interpretation.

\paragraph{Proof.}
Use structural induction on the finite term.  Variables are projections,
zero is constant, and successor is the constant function $1$.  For the
remaining case let
\[
 A(\rho)=\llbracket a\rrbracket_\rho,
 \qquad H(b,\rho)=\llbracket h\rrbracket_{\rho[v\mapsto b]}.
\]
After renaming the bound variable if necessary, the induction hypothesis
says that $A$ is monotone and $H$ is monotone jointly in $b$ and its passive
parameters.  If $\rho\le\sigma$ coordinatewise, then
\[
 H(A(\rho),\rho)\le H(A(\sigma),\sigma).
\]
This is exactly the assertion for $F(n,a;v.h)$.  Dependence on the
counter term may disappear, which causes no problem for monotonicity.
Nested eager binders are handled by the same structural induction.
\hfill$\square$

For completeness, the usual substitution lemma holds for this
interpretation: evaluating a capture-free term substitution is the same
as modifying the corresponding free-variable assignment.  Its proof is
another structural induction.  In the binder case first rename $v$ away
from the substituted variable and replacement term, and then apply the
induction hypotheses to $a$ and $h$.  Consequently the displayed
interpretation respects the formal substitution in both computation
schemas, including arbitrary fixed passive parameters.

\paragraph{Idempotence lemma.}
Every monotone map $H:M\to M$ is idempotent.

\paragraph{Proof.}
The possible pairs $(H(0),H(1))$ are precisely
\[
                       (0,0),\qquad(0,1),\qquad(1,1).
\]
They are respectively the constant-zero map, the identity, and the
constant-one map.  Each satisfies $H(H(b))=H(b)$ for both $b\in M$.
\hfill$\square$

\paragraph{Finite-model theorem.}
The structure $M$ is a model of $T_{\mathrm e}$, and remains a model if
literal-equation induction is replaced by full first-order induction.

\paragraph{Proof.}
Fix arbitrary values of all passive parameters and let $H$ be the
resulting unary interpretation of the update $h$.  By the monotonicity
lemma $H$ is monotone, and therefore idempotent.  If $a^M$ is the value of
the initializer, the zero clause reads
\[
                         H(a^M)=H(a^M).
\]
The left side of the successor clause is again $H(a^M)$, because the
counter is ignored; its right side is $H(H(a^M))$.  These agree by
idempotence.  The substitution lemma justifies these calculations for
the schemas as written, not merely for closed numerical instances.

Coverage holds because the elements are $0$ and $S0=1$.  Also $Sx=1\ne0$
for each $x$.  To check induction, fix any formula $A(x,\vec p)$ and an
assignment of its passive parameters.  If $A(0,\vec p)$ and
\[
             \forall x\,(A(x,\vec p)\to A(Sx,\vec p))
\]
hold, the step at $x=0$ gives $A(1,\vec p)$.  These are all elements of
the domain, so $\forall x\,A(x,\vec p)$ holds.  This argument works for
arbitrary formulas and under any fixed retained context.  The classical
structure is also a one-world intuitionistic Kripke model, so all
intuitionistic first-order rules with equality are valid in it.
\hfill$\square$

\paragraph{The native paired-binder presentation is covered.}
If the eager arithmetic is instead combined with the inherited
Eq/$Q$/$D$ presentation, interpret its clauses by
\[
 Q_yB(y):=B(0)\to\forall y\,B(y),\qquad
 D_{x,y}B(x,y):=\exists x\bigl(B(x,0)\to\forall y\,B(x,y)\bigr).
\]
Hypothesis use, equality transport, $Q$ specialization from its zero
instance, and fresh-witness $D$ elimination are sound in this ordinary
reading.  For case introduction, choose the supplied witness $t$.
The retained base assumption supplies the $y=0$ instance, while the
fresh-successor premise supplies the $y=1=S0$ instance.  For progress
introduction, the retained base assumption and the forward step imply
every instance by the induction just verified.  The usual freshness
conditions keep the witness and passive parameters fixed in these
arguments.  Thus neither removing ordinary connective primitives nor
packaging induction as the native progress rule removes the countermodel.

\paragraph{Corollary: no interpretation of iPRA.}
There is no ordinary finite-dimensional first-order interpretation of
iPRA in $T_{\mathrm e}$, even if its interpreted domain is a definable
subset of $M^d$ modulo a definable equivalence relation.  Consequently
$T_{\mathrm e}$ cannot satisfy the required bi-interpretability condition.

\paragraph{Proof.}
An interpretation would turn every model of $T_{\mathrm e}$ into a model
of iPRA.  In the displayed two-element model, a definable subset of $M^d$
has at most $2^d$ elements, and a quotient cannot increase that number.
A finite collection of interpreted sorts does not change finiteness.
But iPRA proves that all standard successor numerals are distinct, using
successor injectivity and zero--successor separation; hence it has no
finite model.  This is a contradiction.  The same argument permits a
fixed finite tuple of parameters in the interpretation.
\hfill$\square$

\paragraph{Scope and relation to the earlier obstructions.}
This result concerns eager initialization $h(a)$ with an independent
initializer $a$.  It is not the existing almost-unary growth argument
for a tied seed or a zero-seeded restricted iterator.  The supplied
\texttt{local\_scopes} draft also already excludes relational
bi-interpretability for an existential-free tied-seed theory, by finite
witness extraction and finite-profile image counting.  The present
argument obstructs all ordinary finite-dimensional relational
interpretations, including quotients, because it exhibits a finite
model of this different eager theory.  In particular it excludes even
a one-way interpretation of iPRA; it is not a claim to be the first
obstruction to relational bi-interpretability in this research program.
It does not apply to a recursion
update with a separately available counter argument, to the initialized
iterator $I(0,a;h)=a$, or to the successor-seeded iterator
$R(0,a;h)=Sa$.  Those alternatives require their own analyses.
Adding successor injectivity would exclude this countermodel, but would
be an additional arithmetic axiom and would not by itself prove a new
primitive-recursion compiler or bi-interpretation.  This obstruction
rules out the particular apparent simplification; it is not a lower
bound on all iPRA axiomatizations and is not a positive size improvement.
